%%%%%%%%%%%%%%%%%%%%%%%%%%%%%%%%%%%%%%%%%%%%%%%%%%%%%%%%%%%
% --------------------------------------------------------
%  ____  _           
% |  _ \| |__   ___  
% | |_) | '_ \ / _ \ 
% |  _ <| | | | (_) |
% |_| \_\_| |_|\___/ 
%
% LaTeX2e Template
% Version 3.0.0 (24/02/2026)
%
% Author: 
% Guillermo Jimenez (memo.notess1@gmail.com)
% 
% Contributor:
% Eduardo Gracidas (eduardo.gracidas29@gmail.com)
%
% License:
% Creative Commons CC BY 4.0
% --------------------------------------------------------
%%%%%%%%%%%%%%%%%%%%%%%%%%%%%%%%%%%%%%%%%%%%%%%%%%%%%%%%%%%
% --------------------------------------------------------
% Find the latest version on GitHub:
% https://github.com/MemoJimenez/Rho-class
% --------------------------------------------------------
%%%%%%%%%%%%%%%%%%%%%%%%%%%%%%%%%%%%%%%%%%%%%%%%%%%%%%%%%%%

\documentclass[9pt,a4paper,twocolumn,twoside]{rho-class/rho}
\usepackage[english]{babel}
\usepackage{tabularx}
\usepackage{wasysym}
\usepackage{tikz}
\usetikzlibrary{arrows.meta}

%----------------------------------------------------------
% Title
%----------------------------------------------------------

\doctype{Requirement sheet}
\title{Acoustic floor build-up}

%----------------------------------------------------------
% Authors, Affiliations and dates
%----------------------------------------------------------

\author[{\textasteriskcentered},a,1]{Jonathan Bieg}

\affil[a]{jobieg@ethz.ch}

\dates{May 15, 2026}

%----------------------------------------------------------

\theday{\today}

%----------------------------------------------------------
% Abstract and Keywords
%----------------------------------------------------------

\begin{abstract}
    This requirement sheet describes the automated acoustic verification and generation of floor build-ups. The structural cross-section is first evaluated by the mass law. If the acoustic requirements are not fulfilled, a floating screed system consisting of glass wool and cement screed is added. If necessary, gravel is introduced below the floating floor to further improve the acoustic performance. The method distinguishes between normal and increased acoustic requirements and applies an additional planning surcharge for concrete and timber systems. The document also defines the correction terms for gravel layers and insulation inside timber hollow cores and describes the Python classes used for the automated generation of acoustic floor systems.
\end{abstract}

\keywords{Acoustics, Floor Build-up, Mass Law, Floating Screed, Gravel Damping, Timber Floors, Python Automation}

%----------------------------------------------------------

\begin{document}

    \maketitle
    \thispagestyle{firststyle}

%----------------------------------------------------------
\section{Introduction}

The acoustic performance of a floor system depends on the interaction between the structural cross-section and the non-structural floor build-up. Heavy structural systems may already provide sufficient airborne sound insulation, while lighter systems generally require additional layers to satisfy both airborne and impact sound requirements.

\noindent The implemented method generates a suitable floor build-up after the structural cross-section has been defined. The floor build-up therefore adapts to the structural solution instead of being prescribed independently from it. The procedure is intentionally formulated in a small number of robust steps suitable for automated optimisation.

\subsection{Requirements}

The user selects either the requirement level \textit{normal} or \textit{increased}. The corresponding acoustic limit values are:

\begin{table}[H]
\begin{tabular}{l|cc}
    Requirement level & $R_{w,\min}$ & $L_{n,w,\max}$ \\ \hline
    normal & $53\,\mathrm{dB}$ & $51\,\mathrm{dB}$ \\
    increased & $57\,\mathrm{dB}$ & $47\,\mathrm{dB}$
\end{tabular}
\end{table}

\noindent To account for uncertainties during planning and for the additional complexity of the connections, a planning surcharge is applied:

\[
    \Delta_{\mathrm{plan}}=
    \begin{cases}
        5\,\mathrm{dB}, & \text{concrete sections} \\
        8\,\mathrm{dB}, & \text{timber sections}
    \end{cases}
\]

\noindent The verification is therefore carried out with the stricter target values:

\[
    R_{w,\mathrm{target}} = R_{w,\min} + \Delta_{\mathrm{plan}}
\]

\[
    L_{n,w,\mathrm{target}} = L_{n,w,\max} - \Delta_{\mathrm{plan}}
\]

\noindent For the normal requirement level this gives $R_{w,\mathrm{target}}=58\,\mathrm{dB}$ and $L_{n,w,\mathrm{target}}=46\,\mathrm{dB}$ for concrete sections, and $R_{w,\mathrm{target}}=61\,\mathrm{dB}$ and $L_{n,w,\mathrm{target}}=43\,\mathrm{dB}$ for timber sections.

\section{Acoustic calculation}

\subsection{Structural mass}

The structural cross-section is first converted into an area-related mass:

\[
    m'_{\mathrm{section}}=\frac{g_{0,k}}{g}
\]

\noindent where $g_{0,k}$ is the self-weight of the structural cross-section and $g=10\,\mathrm{m/s^2}$ is used in the implementation. If gravel is added, the total base mass becomes:

\[
    m' = m'_{\mathrm{section}} + m'_{\mathrm{gravel}}
\]

\subsection{Mass law}

The airborne sound insulation and the impact sound level of the base construction are estimated from the area-related mass:

\[
    R_w = 37.5 \log_{10}(m') - 42
\]

\[
    L_{n,w} = 164 - 35 \log_{10}(m')
\]

\noindent The structural cross-section and an optional gravel layer are both treated by the mass law.

\subsection{Scope of the model}

The implemented model is intentionally simplified. A more rigorous acoustic prediction would require frequency-dependent input data for the structural systems and the individual floor layers, including measured values for $R(f)$, $L_n(f)$ and the improvements caused by additional layers. Such data is not available consistently for all cross-sections considered in the optimisation. Therefore, a single-number model is used here. It is less detailed than a full spectral calculation and cannot replace a project-specific acoustic verification, but it is transparent, robust, computationally efficient and suitable for comparing many structural alternatives within the automated design workflow.

\subsection{Floating screed system}

If the structural cross-section alone does not fulfil the acoustic target, a floating screed system is added. The floating screed consists of glass wool as spring layer and cement screed as upper mass. The generator checks two predefined stages:

\begin{enumerate}
    \item glass wool with $60\,\mathrm{mm}$ cement screed
    \item glass wool with $85\,\mathrm{mm}$ cement screed
\end{enumerate}

\noindent The floating floor is treated as a mass-spring-mass system. For the airborne sound improvement, the resonance frequency is estimated as:

\[
    f_{0,R}=
    \frac{1}{2\pi}
    \sqrt{
        s'
        \left(
            \frac{1}{m'_{\mathrm{base}}}
            +
            \frac{1}{m'_{\mathrm{screed}}}
        \right)
    }
\]

\noindent with $s'=6\,\mathrm{MN/m^3}$ for the glass wool layer. The improvement of the airborne sound insulation is:

\[
    \Delta R_w =
    74.4
    -
    20 \log_{10}(f_{0,R})
    -
    \frac{R_w}{2}
\]

\noindent For the impact sound improvement, the resonance frequency is calculated as:

\[
    f_{0,L}=160\sqrt{\frac{s'}{m'_{\mathrm{screed}}}}
\]

\noindent and the implemented impact sound reduction of the fast model is obtained from the frequency-band average:

\[
    \Delta L_{n,w}
    =
    \overline{
        30\log_{10}\left(\frac{f}{f_{0,L}}\right)
    }
\]


\subsection{Additional correction terms}

If the floating screed system is still insufficient, gravel is added below the spring layer in discrete steps. The damping effect of the gravel layer is considered by:

\[
    \Delta_1 =
    \min \left[
        \min \left(
            \max \left(
                \frac{m'_{\mathrm{gravel}}}{m'_{\mathrm{section}}},
                1
            \right),
            2
        \right)
        \frac{m'_{\mathrm{gravel}}}{40},
        6
    \right]
\]

\noindent For ribbed timber floors, insulation inside the hollow cores is considered separately from the glass wool layer of the floating screed system. The correction term is:

\[
    \Delta_2 =
    \min \left(
        6 \frac{h_{\mathrm{ins,hollow}}}{0.20},
        7
    \right)
\]

\noindent The correction is only used for \inlinecode{class RibWood}. The hollow-core insulation thickness is taken as the rib height.

\subsection{Verification}

The final acoustic verification is:

\[
    R_w \geq R_{w,\mathrm{target}}
\]

\[
    L_{n,w} \leq L_{n,w,\mathrm{target}}
\]

\noindent The generated floor build-up is accepted as soon as both conditions are fulfilled.

\section{Generated floor build-up}

The build-up is arranged from bottom to top as:

\[
    \text{structural cross-section}
    \rightarrow
    \text{gravel}
    \rightarrow
    \text{glass wool}
    \rightarrow
    \text{cement screed}
    \rightarrow
    \text{parquet}
\]

\noindent The automated workflow is:

\begin{enumerate}
    \item evaluate the structural cross-section by the mass law
    \item if insufficient, add glass wool and $60\,\mathrm{mm}$ cement screed
    \item if insufficient, use glass wool and $85\,\mathrm{mm}$ cement screed
    \item if still insufficient, add gravel below the floating screed system until the acoustic target is reached
\end{enumerate}

\section{Class: AcousticRequirements}

\inlinecode{class AcousticRequirements} stores the selected acoustic requirement level before the section-dependent planning surcharge is added.

\begin{table}[H]
    \begin{tabularx}{0.49\textwidth}{@{}X@{}}
    \hline
    \textbf{\texttt{class AcousticRequirements}} \\ \hline

    \textbf{Attributes} \\
    \texttt{level : string} \\
    \texttt{rw\_min, lnw\_max : double} \\ \hline

    \textbf{Initialization} \\
    \texttt{\_\_init\_\_(level: string = "normal")} \\ \hline
    \end{tabularx}
\end{table}

\section{Class: AcousticResult}

\inlinecode{class AcousticResult} stores the generated floor build-up together with the calculated acoustic values and the contributions of the individual correction terms.

\begin{table}[H]
    \begin{tabularx}{0.49\textwidth}{@{}X@{}}
    \hline
    \textbf{\texttt{class AcousticResult}} \\ \hline

    \textbf{Attributes} \\
    \texttt{floorstruc : FloorStruc} \\
    \texttt{rw, lnw : double} \\
    \texttt{gravel\_thickness, screed\_thickness : double} \\
    \texttt{delta\_gravel, delta\_hollow\_core : double} \\
    \texttt{delta\_rw\_floating, delta\_lnw\_floating : double} \\ \hline
    \end{tabularx}
\end{table}

\section{Class: AcousticFloorGenerator}

\inlinecode{class AcousticFloorGenerator} evaluates the acoustic performance of a floor system and generates the lightest floor build-up that fulfils the section-dependent target values.

\begin{table}[H]
    \begin{tabularx}{0.49\textwidth}{@{}X@{}}
    \hline
    \textbf{\texttt{class AcousticFloorGenerator}} \\ \hline

    \textbf{Methods} \\
    \texttt{section\_mass(section) $\rightarrow$ mass: double} \\
    \texttt{planning\_surcharge(section) $\rightarrow$ surcharge: double} \\
    \texttt{target\_values(section, requirements) $\rightarrow$ rw\_target, lnw\_target: double} \\
    \texttt{hollow\_core\_damping(section) $\rightarrow$ delta\_2: double} \\
    \texttt{gravel\_damping(m\_gravel, m\_section) $\rightarrow$ delta\_1: double} \\
    \texttt{evaluate(section, database\_name, layer\_specs, spring\_stiffness = 6.0) $\rightarrow$ result: AcousticResult} \\
    \texttt{evaluate\_floorstruc(section, floorstruc, spring\_stiffness = 6.0) $\rightarrow$ result: AcousticResult} \\
    \texttt{generate(section, database\_name, requirements = None) $\rightarrow$ result: AcousticResult} \\ \hline
    \end{tabularx}
\end{table}

\printbibliography

%----------------------------------------------------------

\end{document}
